\subsection{The Holt--Klee ceiling: a lift, a doubling, and where the doubling stops}\label{hkd:sec}

\Cref{rem:four-ceilings} leaves $h^{\ast}_{\mathrm{HK}}$ known only at
$m\le5$ and bounded below at $m=6,7$ by the product of
\Cref{lem:hstar-super}, which adds one to the height per dimension. This
subsection records four things: an elementary lift that makes that
``$+1$'' a law rather than a computation on a witness; two explicit
Holt--Klee orientations that push the ceiling at $m=6,7$ to $14$ and
$20$; the observation that both are \emph{doublings}
$2h+2$ realised by the blow-up of \Cref{thm:blowup} with a readout that is
not the parity of the inner height, which is exactly what
\Cref{cor:blowup-parity} did not exclude; and a readout-free obstruction
$(T^{\ast})$ which is the reason the doubling does not iterate. The one-player row of the
table gains the same ``$+1$'' law, by a gadget on the game side.

\begin{lemma}[The facet decomposition of a cube orientation]\label{hkd:facet}
Write the vertices of the $(m+1)$-cube as $(v,b)$ with $v\in\{0,1\}^{m}$,
$b\in\{0,1\}$, the new direction being $m$. For orientations $s_0,s_1$ of
the $m$-cube and $\varphi:\{0,1\}^{m}\to\{0,1\}$ put
\[
  (s_0\!*\!s_1\!*\!\varphi)(v,b)\ :=\ s_b(v)\ \cup\ \{m\ :\ \varphi(v)\ne b\}.
\]
Every orientation of the $(m+1)$-cube is $s_0\!*\!s_1\!*\!\varphi$ for a
unique triple, $s_b$ being its restriction to the facet $\{x_m=b\}$ and
$\varphi(v)=1$ meaning that the edge in direction $m$ at $v$ points away
from the facet $\{x_m=0\}$. Let
\[
 E(s_0,s_1):=\bigl\{\{v,w\}: v\ne w,\ (s_0(v)\oplus s_1(w))\cap(v\oplus w)=\emptyset
   \ \text{ or }\ (s_0(w)\oplus s_1(v))\cap(v\oplus w)=\emptyset\bigr\}.
\]
Then $s_0\!*\!s_1\!*\!\varphi$ is a unique sink orientation if and only if
$s_0$ and $s_1$ are unique sink orientations and $\varphi$ is constant on
every connected component of $E(s_0,s_1)$. If $s_0=s_1$ then
$E(s_0,s_1)=\emptyset$ and every $\varphi$ is admissible. Moreover the
bottom-antipodal step is
\[
  (v,b)\ \longmapsto\ \bigl(v\oplus s_b(v),\ \varphi(v)\bigr),
\]
so the chart used at time $t+1$ is $\varphi$ of the vertex at time $t$.
\end{lemma}

\begin{proof}
Apply the Szab\'o--Welzl condition of \Cref{def:improvement-uso} to each
pair. Two vertices in one facet are separated exactly when that facet's
orientation is a unique sink orientation. For $(v,0)$ and $(v,1)$ the
direction-$m$ bits are $[\varphi(v)\ne0]$ and $[\varphi(v)\ne1]$, which
always differ. For $(v,0)$ and $(w,1)$ with $v\ne w$ the symmetric
difference meets $(v\oplus w)\cup\{m\}$ in the bit $m$ exactly when
$\varphi(v)=\varphi(w)$, and inside $v\oplus w$ exactly when
$(s_0(v)\oplus s_1(w))\cap(v\oplus w)\ne\emptyset$; the pair $(w,0)$,
$(v,1)$ gives the transposed condition. So the only failures are pairs
$\{v,w\}\in E$ with $\varphi(v)\ne\varphi(w)$. When $s_0=s_1=s$ the
Szab\'o--Welzl condition for $s$ makes both intersections nonempty, so $E$
is empty. The step formula: the outer bit of $(v,b)$ is out exactly when
$\varphi(v)\ne b$, and $b\oplus[\varphi(v)\ne b]=\varphi(v)$.
\end{proof}

\begin{remark}\label{hkd:rem-facet}
The decomposition is standard --- it is the routine ``build a cube
orientation from its two facets'' step in the Szab\'o--Welzl and
Schurr--Szab\'o circle of ideas (from memory, unchecked against the
source); \Cref{lem:stack}(a) is its case $s_0=s_1$, $m=2$. It is recorded
here because it makes the search for tall Holt--Klee orientations
\emph{complete}: Holt--Klee is inherited by faces, so both facets of a
Holt--Klee orientation of the $(m+1)$-cube are Holt--Klee orientations of
the $m$-cube, and every Holt--Klee orientation of the $(m+1)$-cube is
$s_0\!*\!s_1\!*\!\varphi$ for such a pair and an admissible $\varphi$. All
the lower bounds below were found in that space.
\end{remark}

\begin{theorem}[The sink lift: $h^{\ast}_{\mathrm{HK}}(m+1)\ge h^{\ast}_{\mathrm{HK}}(m)+1$]\label{hkd:plus-one}
Let $s$ be an acyclic unique sink orientation of the $m$-cube with sink $o$
and bottom-antipodal height $h$, and let
$L(s):=s\!*\!s\!*\!\mathbf 1_{\{o\}}$ be the orientation of the
$(m+1)$-cube with $s$ on both facets and the new edge at $v$ pointing away
from the facet $\{x_m=0\}$ exactly at $v=o$. Then
\begin{enumerate}
\item[(i)] $L(s)$ is an acyclic unique sink orientation with sink $(o,1)$
and bottom-antipodal height exactly $h+1$;
\item[(ii)] $L(s)$ is Holt--Klee if and only if $s$ is.
\end{enumerate}
Consequently $h^{\ast}_{\mathrm{HK}}(m+1)\ge h^{\ast}_{\mathrm{HK}}(m)+1$
and $h^{\ast}(m+1)\ge h^{\ast}(m)+1$ for every $m$.
\end{theorem}

\begin{proof}
$L(s)$ is a unique sink orientation by \Cref{hkd:facet}, the two
facets being equal. Acyclicity: every edge in a direction $<m$ projects to
an edge of $s$ and every edge in direction $m$ projects to a loop, so a
directed cycle projects to a closed directed walk of the acyclic digraph
of $s$, which is constant; the cycle then lies in one fibre
$\{v\}\times\{0,1\}$ and would need $\varphi(v)=1$ and $\varphi(v)=0$. The
sink: $s(o)=\emptyset$ and the direction-$m$ edge at $(o,1)$ is in, since
$\varphi(o)=1=b$.

Height. By \Cref{hkd:facet} the step is
$(v,b)\mapsto(v\oplus s(v),\varphi(v))$, so the first coordinate performs
the bottom-antipodal walk of $s$ whatever $b$ does; from a start with
$h(v_0)=h$ it reaches $o$ after $h$ steps, with $b=\varphi(v_{h-1})=0$
because $v_{h-1}\ne o$, and one further step goes to $(o,1)$, which is the
sink. So the height is $h+1$, and it is not more because the first
coordinate is autonomous.

Holt--Klee. A facet of a Holt--Klee orientation is Holt--Klee, which gives
one direction. For the other, let $F$ be a face of the $m$-cube of
dimension $d\ge1$, so that $F\times\{0,1\}$ is a face of dimension $d+1$;
faces contained in a facet are copies of faces of $s$ and are Holt--Klee
by hypothesis. Let $a$ be the source and $z$ the sink of $s$ on $F$; note
that $o$, having no out-edge at all, is the sink of every face containing
it. Fix $d$ internally vertex-disjoint monotone $a\to z$ paths
$P_1,\dots,P_d$ in $F$, which Holt--Klee supplies; at most one of them is
a single edge, and we index so that $P_1$ is that one if it occurs.

\emph{Case $o\notin F$.} Then $\varphi\equiv0$ on $F$ and every edge in
direction $m$ points from $(v,1)$ to $(v,0)$; the source of
$F\times\{0,1\}$ is $(a,1)$ and its sink is $(z,0)$. Take
\[
\begin{aligned}
 Q_1&:(a,1)\to(a,0)\to[P_1\ \text{below}]\to(z,0), &
 Q_2&:(a,1)\to[P_1\ \text{above}]\to(z,1)\to(z,0),\\
 Q_{i+1}&:(a,1)\to[P_i\ \text{above, as far as } w_i]\to(w_i,0)\to[P_i\ \text{below}]\to(z,0)
   &&(2\le i\le d),
\end{aligned}
\]
where $w_i$ is an interior vertex of $P_i$, which exists because
$P_2,\dots,P_d$ have length at least $2$. These are $d+1$ monotone paths;
$Q_1$ and $Q_2$ use $P_1$ in different copies, $Q_{i+1}$ uses only
interior vertices of $P_i$ in the two copies together with $(w_i,0)$, and
$w_i\notin\{a,z\}$, so they are internally vertex-disjoint.

\emph{Case $o\in F$.} Then $z=o$, $\varphi=\mathbf 1_{\{o\}}$ on $F$, the
edge in direction $m$ points up at $o$ and down elsewhere, the source is
$(a,1)$ and the sink is $(o,1)$. Take the $d$ paths
$(a,1)\to[P_i\ \text{above}]\to(o,1)$ together with
$(a,1)\to(a,0)\to[P_1\ \text{below}]\to(o,0)\to(o,1)$: again $d+1$
monotone, internally vertex-disjoint paths.
\end{proof}

\begin{remark}\label{hkd:rem-plus-one}
Three points. (1) The lift is not a deformed product: a deformed cube has
a uniform last coordinate (\Cref{thm:deformed-flat}), and a uniform
coordinate here means $\varphi$ constant, which gives height $h$, not
$h+1$. The gain comes precisely from making the new coordinate
non-uniform at one vertex. (2) It supersedes the block product of
\Cref{lem:hstar-super} with a $1$-cube in the role of the ``$+1$'' law:
the product preserves Holt--Klee only for some block orders and was
verified on witnesses, whereas $L$ preserves it always. (3) Whether $L(s)$
is an LP orientation whenever $s$ is, is left open; the statement below
for $h^{\ast}_1$ is proved on the game side instead.
\end{remark}

\begin{proposition}[$h^{\ast}_{\mathrm{HK}}(6)\ge14$ and $h^{\ast}_{\mathrm{HK}}(7)\ge20$]\label{hkd:records}
The outmap
\[
\begin{aligned}
 s_6=(&5,10,15,4,9,14,11,8,13,6,7,12,1,0,2,3,\\
      &48,17,51,50,54,61,60,55,30,25,24,59,26,53,52,63,\\
      &32,49,35,34,38,45,44,39,62,57,56,43,58,37,36,47,\\
      &16,33,19,18,22,29,28,23,46,41,40,27,42,21,20,31)
\end{aligned}
\]
of the $6$-cube and
\[
\begin{aligned}
 s_7=(&7,4,21,6,10,11,17,24,29,12,31,14,3,0,25,2,\\
      &23,18,5,20,19,22,8,9,15,28,13,30,27,26,1,16,\\
      &56,113,43,42,102,117,100,103,98,121,32,99,46,127,44,125,\\
      &41,40,54,51,116,101,50,119,48,97,58,59,62,109,60,111,\\
      &120,81,107,106,70,85,68,71,66,89,96,67,110,95,108,93,\\
      &105,104,118,115,84,69,114,87,112,65,122,123,126,77,124,79,\\
      &88,49,75,74,38,53,36,39,34,57,64,35,78,63,76,61,\\
      &73,72,86,83,52,37,82,55,80,33,90,91,94,45,92,47)
\end{aligned}
\]
of the $7$-cube are acyclic unique sink orientations satisfying Holt--Klee
on every face, of bottom-antipodal heights $14$ and $20$; the maximal
walks are
\[
 21,40,22,42,18,33,16,32,0,5,11,7,15,12,13
\]
for $s_6$ (sink $13$; also from $29,53,61$) and
\[
 45,82,36,66,41,80,57,88,40,74,42,10,21,3,5,14,23,30,31,15,13
\]
for $s_7$ (sink $13$; also from $47,109,111$). Hence
\[
  h^{\ast}_{\mathrm{HK}}(6)\ \ge\ 14,\qquad
  h^{\ast}_{\mathrm{HK}}(7)\ \ge\ 20,
\]
and with \Cref{hkd:plus-one},
$h^{\ast}_{\mathrm{HK}}(m)\ge m+13$ for every $m\ge7$. An intermediate
record found by a different method, and kept because its shape is
informative (\Cref{hkd:found}), is the Holt--Klee orientation of the
$7$-cube of height $17$
\[
\begin{aligned}
 (&64,29,31,94,12,51,34,25,8,53,54,55,52,59,42,49,\\
  &28,3,38,9,26,5,24,7,48,15,46,33,36,11,50,13,\\
  &44,45,62,63,40,19,10,41,32,23,2,1,6,27,20,21,\\
  &60,61,14,47,58,43,56,57,22,35,17,16,18,37,4,39,\\
  &0,65,67,30,71,76,93,98,79,72,73,118,75,116,101,106,\\
  &95,92,89,102,115,90,77,88,87,112,117,110,123,100,113,114,\\
  &127,108,125,126,105,104,107,74,103,96,97,66,99,70,85,84,\\
  &109,124,111,78,83,122,121,120,119,86,81,80,91,82,69,68),
\end{aligned}
\]
whose maximal walk is
$68,3,93,57,26,52,14,36,12,56,46,58,43,42,40,8,0,64$ and whose two facets
are Holt--Klee orientations of the $6$-cube of height $13$ with a common
sink.
\end{proposition}

\begin{proof}
Finite computations, made twice in independent implementations (a C engine
and a Python one): the Szab\'o--Welzl condition on all pairs, acyclicity by
depth-first search, the bottom-antipodal heights of all $2^{m}$ starts,
and Holt--Klee by a unit vertex-capacity maximum flow on every face of
dimension at least two ($3^{6}=729$ and $3^{7}=2187$ faces). The
Holt--Klee test was validated on the twelve unique sink orientations of the
$2$-cube, on the $656$ Holt--Klee orientations of the $3$-cube
(\Cref{prop:auso-census}) and on the non-Holt--Klee outmap of
$G^{\sharp}$, where it reproduces the witness $3$-face
$\{x_1,x_2,x_4\}$ at bitmask $4$ of \Cref{cor:seven-two-player}.
\end{proof}

\begin{remark}[How they were found]\label{hkd:found}
$s_7$ was found in the space of \Cref{hkd:facet}: its two facets are
Holt--Klee orientations of the $6$-cube of height $13$ with the same sink
(they are not equal), and $\varphi=\mathbf 1_{\{o,3\}}$ has weight two. Its walk
spends three steps changing chart --- $(4,1)\to(3,0)\to(29,1)\to(57,0)$
--- and then runs the whole height-$13$ walk of the lower facet and closes
with the step of \Cref{hkd:plus-one}, $17=3+13+1$. The pattern is the
same at the previous level: a Holt--Klee orientation of the $6$-cube of
height $13$ has two isomorphic height-$11$ facets with a common sink and
$\varphi$ of co-weight two, and $13=1+11+1$. So the gain over
\Cref{hkd:plus-one} is the length of a chart-changing preamble, and
bounding that is the same problem one dimension down. As a control, over
$30\,000$ random pairs taken from the automorphism orbit of the height-$11$
witness of \Cref{prop:hkfive}, with every admissible $\varphi$
enumerated, orientations of the $6$-cube of bottom-antipodal height up to
$21$ occur, but every one of height at least $16$ is cyclic, every acyclic
one of height $14$ or $15$ fails Holt--Klee, and the Holt--Klee ones stop
at $13$; the height-$14$ orientation $s_6$ comes from the blow-up instead
(\Cref{hkd:doubling}). A caution on method that cost us the first
search: a tall Holt--Klee orientation need not have tall facets. All
twelve facets of $s_6$ have height $6$, $7$ or $8$, against the ceiling
$h^{\ast}_{\mathrm{HK}}(5)=11$, so restricting a facet search to
height-$11$ facets --- which is what produced the height-$13$ examples ---
misses $s_6$ entirely.
\end{remark}

\begin{proposition}[The generalised blow-up: an arbitrary readout]\label{hkd:readout}
For an orientation $s$ of the $m$-cube, $z\in\{0,1\}^{m}$ and a
\emph{readout} $\varphi:\{0,1\}^{m}\to\{0,1\}$ let $B_\varphi(s,z)$ be the
outmap of the $(m+2)$-cube given by the table of \Cref{thm:blowup} with
$\varphi(v)$ in place of the parity of $h(v)$. Then:
\begin{enumerate}
\item[(a)] If $s$ is an acyclic unique sink orientation then so is
$B_\varphi(s,z)$, for \emph{every} $z$ and every $\varphi$.
\item[(b)] If $\varphi$ alternates along a bottom-antipodal walk of $s$
from a vertex $u$ of height $h$ to the sink $o$, and $z=o\oplus u$, then
$B_\varphi(s,z)$ has height $2h+2$.
\item[(c)] $B_\varphi(s,z)$ is Holt--Klee only if, for every $2$-face $F$
of $s$, $\varphi$ is unbalanced on $F$ or $\varphi(\mathrm{src}F)\ne
\varphi(\mathrm{snk}F)$ \emph{(P$_\varphi$)}, and the translation
condition \emph{(T)} of \Cref{lem:blowup-faces} holds; \emph{(T)} does not
involve $\varphi$.
\end{enumerate}
\end{proposition}

\begin{proof}
(a) The Szab\'o--Welzl condition. Two vertices in one layer are separated
by the layer's inner orientation. For two vertices in different layers with
the same inner $v$, the outer outmaps of \Cref{thm:blowup} differ on the
layer difference in all eight cases. For different layers and different
inner vertices, the same check leaves only the pairs of layers
$\{(1,0),(0,1)\}$, $\{(1,0),(1,1)\}$ and $\{(0,1),(1,1)\}$ at one readout
combination each, and there both layers carry $s$, which separates them.
Acyclicity: the inner coordinate moves only along edges of the acyclic
$s$ or of the acyclic $s(\cdot\oplus z)$, so any directed cycle has
constant inner vertex and lies in the $2$-face of layers, whose
orientation has the sink $(0,0)$ and is not a directed cycle.
(b) The layer dynamics of the table sends $(1,0)\mapsto(0,1)$ at
$\varphi=0$ and $(0,1)\mapsto(1,0)$ at $\varphi=1$, and drops to $(0,0)$
in the two remaining cases; so an alternating readout keeps the walk in
the layers $(1,0),(0,1)$ for the $h$ steps of the inner walk and for one
more, then enters $(0,0)$, where the inner map is $s(\cdot\oplus z)$ and
the inner vertex $o$ reads as $o\oplus z=u$, so the inner walk runs again,
$h$ steps. Total $2h+2$.
(c) The $3$-faces of $B_\varphi$ listed in \Cref{lem:blowup-faces} are
stacks with $\varphi$, with $1-\varphi$, and (twice) with the constant
$0$; \Cref{lem:stack}, which we re-verified here over all $192$ and all
$144$ pairs, gives the two conditions, and Holt--Klee is inherited by
faces.
\end{proof}

\begin{proposition}[The readout-free obstruction $(T^{\ast})$]\label{hkd:tstar}
For an orientation $s$ of the $m$-cube and $z\in\{0,1\}^{m}$ let
\[
  D(s,z)\ :=\ \bigl(s(\cdot\oplus z)\bigr)\!*\!s\!*\!\mathbf 0
\]
be the orientation of the $(m+1)$-cube with $s(\cdot\oplus z)$ on the
facet $\{x_m=0\}$, $s$ on the facet $\{x_m=1\}$ and every edge in
direction $m$ pointing towards the facet $\{x_m=0\}$. If $B_\varphi(s,z)$
is Holt--Klee for even one readout $\varphi$, then $D(s,z)$ is Holt--Klee.
Restricted to $2$-faces this is exactly condition \emph{(T)} of
\Cref{lem:blowup-faces}; at higher dimensions it is strictly stronger.
\end{proposition}

\begin{proof}
By the table of \Cref{thm:blowup} the outer coordinate $\alpha$ lies in
$O_{10}$ for both readout values and in no $O_{00}$, so every
$\alpha$-edge between the layers $(1,0)$ and $(0,0)$ points towards
$(0,0)$, whatever $\varphi$ is; and the inner parts on those two layers
are $s$ and $s(\cdot\oplus z)$. Hence for every face $F$ of the $m$-cube
the face of $B_\varphi(s,z)$ on the directions of $F$ together with
$\alpha$, taken at $\beta=0$, is isomorphic to the corresponding face of
$D(s,z)$; Holt--Klee is inherited by faces. The $2$-face case is
\Cref{lem:stack}(b), which is \emph{(T)}. Strictness: for the seed
$s$ of \Cref{hkd:doubling}(d) below at $m=5$ with $z=21$ all $80$
$2$-faces satisfy \emph{(T)} while $D(s,21)$ fails Holt--Klee already on a
$4$-face, and correspondingly no readout makes $B_\varphi(s,21)$
Holt--Klee: the failing faces are the two $4$-faces on
$\{x_0,x_3,x_4,\alpha\}$ and $\{x_0,x_3,x_4,\beta\}$ at the origin and the
$5$-face containing them, for every readout that satisfies
\emph{(P$_\varphi$)}.
\end{proof}

\begin{remark}\label{hkd:rem-tstar}
$(T^{\ast})$ is one Holt--Klee test on a cube of dimension $m+1$, so it is
a cheap filter, and it is the filter that produced $s_7$: of $85\,671$
Holt--Klee orientations of the $5$-cube generated by the facet search,
$6\,274$ pass \emph{(T)} for some $z=o\oplus u$ but only $188$ pass
$(T^{\ast})$ with $h(u)\ge6$, and of those only three have $h(u)=9$;
those three are exactly where a readout was found.
\end{remark}

\begin{proposition}[The blow-up doubles once, with a non-parity readout]\label{hkd:doubling}
\begin{enumerate}
\item[(a)] Over all $728$ acyclic unique sink orientations $s$ of the
$3$-cube, all $8$ translations $z$ and all $256$ readouts $\varphi$
--- $1\,490\,944$ orientations of the $5$-cube --- the greatest
bottom-antipodal height of a Holt--Klee $B_\varphi(s,z)$ is $8$; the
doubled value $2h^{\ast}(3)+2=10$ is not attained, nor is $9$.
\item[(b)] At $m=4$ it is attained. For the Holt--Klee seed
\[
  s=(0,1,3,2,6,13,12,7,14,9,8,11,10,5,4,15)
\]
of the $4$-cube, of height $6$ with sink $o=0$ and maximal vertices
$\{5,13\}$, and $z=o\oplus13=13$, the readout $\varphi$ with support
$\{1,8,9,10,12\}$ makes $B_\varphi(s,z)$ a Holt--Klee acyclic unique sink
orientation of the $6$-cube of height $14=2\cdot6+2=
2h^{\ast}_{\mathrm{HK}}(4)+2$; it is the orientation $s_6$ of
\Cref{hkd:records}. Taking $\varphi$ to be the parity of the inner
height and flipping it at the \emph{other} maximal vertex $5$ also works.
So the obstruction of \Cref{cor:blowup-parity} is an obstruction to the
parity readout only.
\item[(c)] At $m=5$ it is attained again, at height $9$ rather than at
the ceiling $11$. For the Holt--Klee seed
\[
 s^{(5)}=(24,17,11,10,6,21,4,7,2,25,0,3,14,31,12,29,
          9,8,22,19,20,5,18,23,16,1,26,27,30,13,28,15)
\]
of the $5$-cube, of height $9$ with sink $o=10$, and $u=13$ (of height
$9$), $z=o\oplus u=7$, the readout $\varphi$ with support
\[
 \{0,2,3,10,12,14,16,17,18,19,22,24,26,27,28,30\}
\]
makes $B_\varphi(s^{(5)},7)$ a Holt--Klee acyclic unique sink orientation
of the $7$-cube of height $20=2\cdot9+2$; it is the orientation $s_7$ of
\Cref{hkd:records}. Three seeds pass $(T^{\ast})$ with $h(u)=9$ and
between them yield $96$ such readouts.
\item[(d)] The two conditions are of different strength, and only
$(T^{\ast})$ is predictive. Of the $785$ Holt--Klee classes of the
$4$-cube of height $5$ or $6$, exactly $506$ admit a $z=o\oplus u$
satisfying \emph{(T)} and $4$ admit one with $h(u)=6$; enumerating all
$2^{16}$ readouts over those four gives exactly $36$ Holt--Klee blow-ups
of height $14$. At $m=5$, of $85\,671$ Holt--Klee orientations produced by
the facet search, $6\,274$ pass \emph{(T)} for some $z=o\oplus u$, with
$h(u)$ up to $9$; the seed
\[
 (8,25,7,2,5,6,3,4,1,0,11,10,14,15,13,12,
  16,9,19,18,22,23,28,21,24,17,27,26,30,29,20,31)
\]
with $z=21$, $u=28$, $h(u)=8$ passes \emph{(T)} on all $80$ of its
$2$-faces, satisfies \emph{(P$_\varphi$)} for a readout, and yet admits
no Holt--Klee blow-up at all: $D(s,21)$ fails Holt--Klee on the $4$-face
$\{x_0,x_3,x_4,x_5\}$ at the origin, and every \emph{(P)}-feasible readout
leaves a defect of exactly three faces of $B_\varphi(s,21)$, the two
$4$-faces on $\{x_0,x_3,x_4,\alpha\}$ and $\{x_0,x_3,x_4,\beta\}$ at the
origin and the $5$-face containing them. Only $188$ of the $85\,671$ pass
$(T^{\ast})$ with $h(u)\ge6$, and just $3$ with $h(u)=9$.
\item[(e)] The chain stops there on our material. None of the $36$
height-$14$ orientations of the $6$-cube of (b), none of $502$ Holt--Klee
orientations of the $6$-cube of height $13$, and none of the $96$
height-$20$ orientations of the $7$-cube of (c) admits a translation
passing $(T^{\ast})$ with $h(u)\ge9$ (resp.\ $\ge14$); so none of them can
be blown up again, whatever the readout.
\end{enumerate}
\end{proposition}

\begin{proof}
Exhaustive computations with the verified engine of
\Cref{hkd:records}; (a) is a complete scan of the stated space;
\emph{(T)} is tested on all $\binom m2 2^{m-2}$ $2$-faces and
$(T^{\ast})$ by one Holt--Klee test on the $(m+1)$-cube. In (b) and (c)
the readouts were checked from the definition of $B_\varphi$ and the
resulting orientations independently re-verified in a second
implementation.
\end{proof}

\begin{remark}\label{hkd:rem-doubling}
\Cref{hkd:doubling}(b) answers, at the first size where it can be
asked, the question raised after \Cref{cor:blowup-transl}: a
dimension-raising doubling \emph{can} preserve Holt--Klee, so
$h^{\ast}_{\mathrm{HK}}(6)\ge2h^{\ast}_{\mathrm{HK}}(4)+2$; and (c) does
it once more, from a seed of height $9$ rather than of the ceiling
height, which is why $20$ and not $24$. What (a) and (e) say is that this
is still not a family. At $m=3$ the rule cannot double at all, for any
readout. Above that the binding constraint is the readout-free
$(T^{\ast})$ of \Cref{hkd:tstar}, which is a Holt--Klee condition on
a cube one dimension up and therefore gets harder as $m$ grows: at $m=5$
it holds for three seeds in $85\,671$ at $h(u)=9$ and for none at $h(u)=10$
or $11$, and at $m=6,7$ for none of the tall orientations we have. So
\[
  h^{\ast}_{\mathrm{HK}}(m+2)\ \ge\ 2h^{\ast}_{\mathrm{HK}}(m)+2
\]
is now verified at $m=4$ ($14\ge14$), false for this rule at $m=3$, and
open from $m=5$ on --- at $m=5$ the rule delivers $20$ against the
$24$ the inequality asks for. What a superpolynomial family would need is
a construction that supplies its own $(T^{\ast})$-clean translation at
every level; the blow-up does not, and nothing here decides whether
$h^{\ast}_{\mathrm{HK}}$ is superpolynomial.
\end{remark}

\begin{theorem}[One player: $h^{\ast}_{1}(m+1)\ge h^{\ast}_{1}(m)+1$]\label{hkd:one-plus-one}
Let $G$ be a stopping SSG with $\Vmin=\emptyset$, $|\Vmax|=m$, and an
all-switches run $\sigma_0,\dots,\sigma_L$ with $L\ge1$. Then there is a
stopping SSG $G'$ with $\Vmin=\emptyset$ and $|\Vmax|=m+1$ whose
all-switches run from an explicit start has length $L+1$; if $G$ is
nondegenerate then $G'$ can be taken nondegenerate. Hence
$h^{\ast}_{1}(m+1)\ge h^{\ast}_{1}(m)+1$ and, by
\Cref{prop:oneplayer-lp}, $h^{\ast}_{\mathrm{LP}}(m+1)\ge
h^{\ast}_{\mathrm{LP}}(m)+1$ along the same witnesses. With
$h^{\ast}_{1}(6)\ge12$ of \Cref{prop:oneplayer-runs},
\[
  h^{\ast}_{1}(m)\ \ge\ m+6\qquad(m\ge6),
\]
so in particular $h^{\ast}_{1}(7)\ge13$, against the $\ge12$ of
\Cref{rem:four-ceilings}.
\end{theorem}

\begin{proof}
Work in the harmonic normal form of \Cref{prop:oneplayer-runs}. The value
vectors of a run are pairwise distinct and nondecreasing
(\Cref{lem:switch}, \Cref{lem:flat-class}), so some $x_j\in\Vmax$ has
$\val_{\sigma_{L-1}}(x_j)<\val_{\sigma_L}(x_j)$. Fix
$\varepsilon:=2^{-D_0}$ with $D_0$ large enough that
$(1-\varepsilon)\val_{\sigma_{L-1}}(x_j)<(1-\varepsilon)\val_{\sigma_L}(x_j)$
still holds, and choose a dyadic $\theta$ strictly between them; by
\Cref{lem:denominator-sharp} two distinct values of $G$ differ by at least
$2^{-2a}$, so $\theta$ of denominator $2^{2a+2}$ exists. Adjoin one Max
vertex $y$ with the two rows
\[
  (y,0):\ p=0,\ q=\theta,\qquad (y,1):\ p=(1-\varepsilon)e_j,\ q=0,
\]
and leave every old row unchanged, with a zero entry in the new
coordinate; by \Cref{lem:dyadic-row} the two rows are realised by chains
of average vertices, and every row still leaks, so $G'$ is stopping. No
old vertex sees $y$, so $\val_\sigma$ on $\Vmax(G)$ and the switchable
sets of the old vertices are unchanged. Start at $\sigma_0$ with
$\sigma_0(y)=0$. For $t<L$ the alternative at $y$ has value
$(1-\varepsilon)\val_{\sigma_t}(x_j)\le
(1-\varepsilon)\val_{\sigma_{L-1}}(x_j)<\theta$, so $y$ is not switchable
and the run is that of $G$; at $\sigma_L$ the old switchable set is empty
while the alternative at $y$ has value
$(1-\varepsilon)\val_{\sigma_L}(x_j)>\theta$, so $S_{\sigma_L}=\{y\}$ and
the run continues one round. Nondegeneracy: $\theta$ ranges over an
interval, so it may be chosen to avoid the finitely many numbers
$(1-\varepsilon)\val_\sigma(x_j)$.
\end{proof}

\begin{remark}\label{hkd:rem-one-plus-one}
The gadget is the one \Cref{cor:hstar-one} uses to show $h^{\ast}_1$
nondecreasing, with one change that is the whole point: the new Max vertex
is aimed at a vertex \emph{of the game} rather than at a fresh coin, and
started at the constant, so that it wakes at the \emph{end} of the run and
not at round zero. Iterating it five times on the $m=6$ game of
\Cref{prop:oneplayer-runs} (start $25$, run length $12$) gives explicit
nondegenerate stopping one-player games with
\[
  m=7,8,9,10,11 \quad\text{and run lengths}\quad 13,14,15,16,17,
\]
their runs being $25,14,33,6,37,46,45,44,16,56,24,28,12,76,204,460,972,
1996$ truncated to the right length; we verified each from its stored
normal form in exact rational arithmetic, together with stoppingness and
the absence of any tied incidence over all $2^{m}$ strategies. The
thresholds are $\theta=985727/2^{20}$, $252284505/2^{28}$,
$504445825/2^{29}$, $4034581355/2^{32}$, $2016798175/2^{31}$, so the
denominators grow slowly and the games stay small --- at $m=11$,
\Cref{lem:dyadic-row} bounds the size by $N\le717$ --- but the proof gives
only $D\le2a+2$ per step, hence no polynomial bound on the size of the
$k$-th iterate; the measured growth is linear.
\end{remark}

The four ceilings now read
\[
\begin{array}{l|ccccccc}
m & 1 & 2 & 3 & 4 & 5 & 6 & 7\\ \hline
h^{\ast}(m) & 1 & 2 & 4 & 7 & 12 & \ge16 & \ge26\\
h^{\ast}_{\mathrm{HK}}(m) & 1 & 2 & 4 & 6 & 11 & \ge14 & \ge20\\
h^{\ast}_{\mathrm{LP}}(m) & 1 & 2 & 4 & 6 & \ge9 & \ge12 & \ge13\\
h^{\ast}_{1}(m) & 1 & 2 & 4 & 6 & \ge9 & \ge12 & \ge13
\end{array}
\]
with $h^{\ast}_{\mathrm{HK}}(m)\ge m+13$ for $m\ge7$
(\Cref{hkd:plus-one}, \Cref{hkd:records}) and
$h^{\ast}_{1}(m),h^{\ast}_{\mathrm{LP}}(m)\ge m+6$ for $m\ge6$
(\Cref{hkd:one-plus-one}). These are the first unbounded lower bounds
in the last three rows; they are linear, and whether any of the three
grows superpolynomially is untouched.